%% ----------------------------------------------------------------
%% Chapter_2.tex
%% ----------------------------------------------------------------
\chapter{Literature Survey}
\label{Chapter:Literature}

\section{Existing Systems and Studies}
Several studies and institutional practices have explored the use of feedback systems to evaluate academic quality and student satisfaction. Traditional feedback mechanisms primarily focus on parameters such as teaching effectiveness, infrastructure, and course delivery. In addition, some academic studies analyze student performance using quantitative indicators such as examination scores, attendance records, and internal assessments.

With the emergence of education analytics, data-driven approaches have been adopted to support academic decision-making. However, most existing systems rely predominantly on structured numerical data and provide limited support for analyzing qualitative information, such as alumni career experiences, industry exposure, and real-world application of academic knowledge.

Employability tracking systems, where available, are often manual, fragmented, and loosely connected to curriculum design processes, thereby limiting their effectiveness in guiding curriculum improvement.

\section{State of the Art}
Contemporary education analytics platforms provide dashboard-based reporting and summary-level insights to assist institutions in monitoring academic performance and learning outcomes. Some advanced systems incorporate data mining and basic machine learning techniques; however, their application remains largely confined to academic metrics such as grades, progression rates, and assessment outcomes.

Integration of alumni career data with curriculum analysis remains limited in current solutions. Most state-of-the-art systems lack a structured methodology to evaluate the alignment between academic subjects and industry-required skills. Consequently, curriculum revision decisions are often based on indirect indicators rather than concrete employability evidence. The proposed system, EduFeedback AI, aims to extend the state of the art by introducing an integrated framework that combines feedback analytics, employability data, and curriculum relevance assessment.

\section{Market Research}
Recent trends in higher education and workforce development indicate a growing disconnect between academic curricula and evolving industry requirements. Market research from global and national agencies highlights the urgency for data-driven curriculum evaluation and employability analytics.

According to the World Economic Forum (Future of Jobs Report), skill disruption across industries remains significant. Employers estimate that a large portion of employees’ core skills will change within a short time frame due to technological advancements, automation, and artificial intelligence. The report indicates that nearly 40\% of core job skills are expected to change by 2030, emphasizing the need for continuous reskilling and curriculum adaptation.

In the Indian context, employability-driven curriculum reform has become even more critical. The India Skills Report 2025 indicates a rapid increase in demand for industry-ready professionals, particularly in emerging technology domains such as artificial intelligence, data analytics, and automation. The report estimates that India had approximately 4.16 lakh AI-ready professionals in 2023, with projections reaching nearly 10 lakh professionals by 2026. A national employability index shows the capability of graduates to secure jobs:

\begin{itemize}
    \item India’s employability rate in 2025-2026: 56.35\%
\end{itemize}

This means slightly over half of graduates are considered ready for employment according to market expectations.

\begin{figure}[!h]
    \centering
    % \includegraphics[width=0.7\linewidth]{figures/employability_rate.png}
    \caption{India’s employability rate}
    \label{fig:employability_rate}
\end{figure}

A comparison of employability rates across different engineering domains reveals noticeable variation, as illustrated in the bar graph below.

\begin{figure}[!h]
    \centering
    % \includegraphics[width=0.7\linewidth]{figures/employability_rate_sectors.png}
    \caption{India’s employability rate across different sectors}
    \label{fig:employability_sectors}
\end{figure}

Market analysis further reveals that accreditation frameworks such as NAAC and NBA increasingly emphasize outcome-based education, graduate employability, and stakeholder feedback. Institutions are required to provide measurable evidence demonstrating how curricula contribute to career readiness and long-term professional success.

In this context, the need for a system like EduFeedback AI becomes evident. By systematically collecting alumni career data, faculty input, and curriculum feedback, and by analyzing this information using data-driven techniques, the system supports evidence-based curriculum improvement. This approach helps institutions respond proactively to market demands, reduce skill gaps, and improve graduate employability outcomes.

\section{Research Gap and Limitations of Existing Systems}
The literature survey identifies several limitations in existing approaches:
\begin{itemize}
    \item Absence of continuous alumni career tracking mechanisms
    \item Limited analysis of unstructured feedback data
    \item Lack of direct mapping between curriculum components and industry skill requirements
    \item Insufficient analytical support for data-driven curriculum revision
\end{itemize}

These gaps highlight the necessity for an integrated framework capable of combining feedback, employability analysis, and curriculum evaluation. EduFeedback AI is proposed to address these shortcomings through a unified and extensible system architecture.

\section{Motivation for Project Selection}
The selection of this project is motivated by the growing disconnect between academic curricula and industry expectations observed across higher education institutions. While feedback from alumni is frequently collected, it is rarely analyzed in a structured and actionable manner.

To address this challenge, the project seeks to answer the following key questions, which form the core motivation of the proposed work:
\begin{itemize}
    \item Which academic subjects and skills are most relevant to current industry job roles?
    \item What skills are graduates acquiring after graduation that were not adequately covered during their academic program?
    \item How effectively does the existing curriculum prepare students for real-world employment?
    \item What skill gaps exist between academic training and industry expectations across different domains?
    \item How can alumni feedback be converted into measurable indicators to support data-driven curriculum revision?
\end{itemize}

EduFeedback AI seeks to address this gap by providing a scalable and data-driven solution that transforms feedback into measurable indicators of curriculum relevance and employability. The project also aligns with contemporary research directions in education analytics and institutional quality assurance, making it suitable for academic and practical deployment.
